\section{Vorgehen und Ablauf}
\label{sec:ablauf}

\subsection{Zeitplan}
Gleich zu Beginn Ihrer Arbeit sollten Sie eine grobe Zeitplanung, z. B. in der Form eines Balkenplanes und eines Meilensteinplans, anlegen und diese während der Arbeit ständig korrigieren und fortschreiben. Durch solche einfachen Mittel des Projektmanagements können Sie unnötigen Zeitverlust vermeiden, der oft dadurch entsteht, dass man sich mit Randproblemen zu lange aufhält. Spreche diese Zeitpläne mit deiner Betreuerin/ deinem Betreuer durch.

Die Zeit von der Ausgabe des Themas bis zur Ablieferung der Ausarbeitung der Abschlussarbeit darf bei Bachelorarbeiten zwölf Wochen, bei Masterarbeiten XX Wochen nicht überschreiten. Auf Antrag des Kandidaten kann der Vorsitzende des Prüfungs-ausschusses die Bearbeitungszeit in Ausnahmefällen in Einverständnis mit dem Aufgabensteller auf maximal das doppelte verlängern. Wird die Arbeit nicht fristgemäß eingereicht, so gilt sie als abgelegt und nicht bestanden.

Die Erfahrung zeigt, dass für die Abfassung der Dokumentation ca. drei Wochen einzuplanen sind.



Nach Beendigung der Arbeit sind alle Schlüssel, entliehenen Geräte, überzähligen Bauteile, Software, etc., abzugeben. Die Bestätigung der Rückgaben erfolgt auf dem Material- und Rückgabeschein, die Löschung der Accounts auf der Rechnerberechtigung. Beide sind ausgefüllt bei Frau Werner (Zi. N2506) abzugeben. 


